%% kapa: exam=1 round=27 year=2016 period=1 form=A subject=civil_law subject_ko=민법 seq=40 number=40
%% source: Q-Net 감정평가사 기출문제 (cst00309) artlSeq=5169402 file=extracted/1008/5169402/제27회 감정평가사 1차 문제/제27회 감정평가사 1교시문제(A형).pdf
%% conversion: kordoc 4.24.51 -> tools/gen1.py (auto); review=proofread-sonnet; flags=none
\begin{question}{40}
전세권에 관한 설명으로 옳지 않은 것은? (다툼이 있으면 판례에 따름)
\begin{choices}
  \item 전세권자가 소유자의 승낙 없이 전세권을 제3자에게 양도한 점만으로는 전세권에 대한 소멸청구사유가 되지 않는다.
  \item 타인의 토지에 있는 건물에 설정된 전세권의 효력은 그 건물의 소유를 목적으로 한 토지임차권에도 미친다.
  \item 전세권자는 통상의 필요비와 유익비를 지출한 경우, 전세권설정자에게 그 상환을 청구할 수 있다.
  \item 전세권의 목적물의 일부가 불가항력으로 인하여 멸실된 때에는 그 멸실된 부분의 전세권은 소멸한다.
  \item 건물의 일부에 대하여만 전세권이 설정되어 있는 경우에 그 전세권자는 건물 전부의 경매를 청구할 수 없다.
\end{choices}
\end{question}
\answer{3}
